% PAGE-BUDGET TARGET: 8 pages
% HARD MAX: 9 pages
% BLUEPRINT PHASE 4: Draft Sections 7--8 and Appendix B.
\section{Endpoint packet: reconstruction, non-triviality, and local-net universality}\label{sec:endpoint}
\budgetnote{8}{9}

The preceding sections establish the sharp-local state and the continuum vacuum gap. This section packages the claim-facing consequences: the Euclidean OS dossier on the full sharp-local algebra, the Wightman--Haag--Kastler reconstruction, field correspondence for the renormalized local gauge-invariant operators produced by Lane~A, a non-triviality witness, and the faithful-Wilson universality theorem that identifies the exact local-net endpoint of the paper.

\subsection{Reconstruction and non-triviality packet}

\begin{theorem}[Euclidean OS dossier on the full sharp-local algebra]\label{thm:euclidean-os-dossier}
Let $\sharpstate$ be the full sharp-local state of Corollary~\ref{cor:full-inductive-union}. Then the Schwinger distributions generated by the renormalized local fields of Lane~A satisfy the following properties on the full sharp-local algebra $\Aloc$:
\begin{enumerate}[label=\textup{(\roman*)},leftmargin=2.25em]
\item Euclidean invariance under translations and rotations.
\item Bosonic symmetry under permutation of insertions.
\item Osterwalder--Schrader reflection positivity on the full positive-time sharp-local algebra.
\item Regularity and temperedness of the Schwinger distributions.
\item OS reconstruction on the full positive-time sharp-local algebra, yielding a Hilbert space
\[
(\mathcal H_{\mathrm{loc}},\Omega_{\mathrm{loc}},H_{\mathrm{loc}}),
\]
with $\Omega_{\mathrm{loc}}=[1]$ cyclic for the full sharp-local positive-time algebra. Moreover the bounded positive-time base algebra $A_+$ embeds densely in $\mathcal H_{\mathrm{loc}}$.
\end{enumerate}
\end{theorem}

\proofsynopsis{Items~(i)--(iii) are the Euclidean invariance, symmetry, and reflection-positivity clauses already verified capwise in the Lane~A finite-cap theorem and then passed to the inductive union. Temperedness is imported from the sharp-local Euclidean control package, and the OS reconstruction in item~(v) is the standard reconstruction theorem applied to the full sharp-local Schwinger family. The only theorem-facing point retained in the core paper is that the constructive Lane~A output is already strong enough to trigger the standard OS machinery on the full sharp-local algebra.}
\proofhome{Companion III, Section~1.}

\begin{theorem}[Wightman--Haag--Kastler reconstruction]\label{thm:wightman-reconstruction}
Let $\sharpstate$ be as above. Then the Euclidean Schwinger functions of the full sharp-local theory admit a Minkowski reconstruction with the following properties:
\begin{enumerate}[label=\textup{(\roman*)},leftmargin=2.25em]
\item there exists a Hilbert space $\mathcal H_{\mathrm{loc}}$, a vacuum vector $\Omega_{\mathrm{loc}}$, and a strongly continuous unitary representation $U$ of the Poincar\'e group whose joint spectrum lies in the closed forward light cone;
\item for every renormalized local gauge-invariant field generated by Lane~A there is a corresponding operator-valued tempered distribution on Minkowski space whose Wightman distributions are the analytic continuations of the Euclidean Schwinger distributions;
\item these fields satisfy locality and generate the vacuum Haag--Kastler net of the local gauge-invariant sharp-local theory;
\item the reconstruction is unique up to unitary equivalence once the Euclidean Schwinger functions are fixed.
\end{enumerate}
\end{theorem}

\proofsynopsis{Once \cref{thm:euclidean-os-dossier} is in place, the theorem is the standard OS analytic-continuation/reconstruction package applied to the full sharp-local Schwinger family. The core paper keeps the endpoint statement only; the full reconstruction proof and imported uniqueness clause are sent to Companion~III.}
\proofhome{Companion III, Section~2.}

\begin{corollary}[Field correspondence]\label{cor:field-correspondence}
For every exact-dimension quotient label $[P]$ and every test function $f$, let
\[
O^{\mathrm{ren}}_{[P]}(f;\mu)
\]
denote the corresponding renormalized sharp-local field constructed by Lane~A. Then in the Wightman reconstruction of \cref{thm:wightman-reconstruction} there exists a unique operator-valued tempered distribution $\Phi_{[P]}$ whose Euclidean Schwinger distributions are the analytic continuations of the mixed Schwinger functions of $O^{\mathrm{ren}}_{[P]}(f;\mu)$. This correspondence is compatible with finite linear combinations, derivatives, quotient relations, finite engineering truncations, and the inductive-union passage to the full sharp-local algebra.

Consequently the reconstructed local field content is exactly the family of gauge-invariant local differential polynomials in the curvature and its covariant derivatives represented by the Lane~A sharp-local algebra.
\end{corollary}

\begin{proof}
This is the field-correspondence output of \cref{thm:wightman-reconstruction} applied to the canonical exact-dimension quotient blocks and the inductive-union sharp-local algebra already constructed in Section~\ref{sec:laneA}.
\end{proof}

\begin{corollary}[Minkowski Hamiltonian inherits the gap]\label{cor:minkowski-gap}
Let $P^0$ be the Minkowski Hamiltonian in the Wightman reconstruction of \cref{thm:wightman-reconstruction}. Assume the continuum sharp-local vacuum-gap theorem Corollary~\ref{cor:continuum-vacuum-gap}. Then
\[
\operatorname{Spec}(P^0)\subset \{0\}\cup[m_*,\infty).
\]
In particular, the reconstructed local gauge-invariant sharp-local theory has positive energy, a unique vacuum, and a strictly positive Minkowski mass gap in its vacuum sector.
\end{corollary}

\begin{proof}
Under OS reconstruction, Euclidean time translations analytically continue to the Minkowski time-translation group. The Minkowski Hamiltonian therefore coincides with the OS generator on the reconstructed Hilbert space, so the spectral inclusion transfers directly from Corollary~\ref{cor:continuum-vacuum-gap}.
\end{proof}

\begin{theorem}[Non-triviality witness]\label{thm:nontriviality}
Assume the full sharp-local OS/Wightman construction of \cref{thm:wightman-reconstruction} together with the Route~1 vacuum-gap theorem Corollary~\ref{cor:continuum-vacuum-gap}. Let $E^{\mathrm{ren}}$ be the canonical centered CP-even dimension-4 scalar witness produced by the Lane~A normalization scheme, and let $\varphi_{*,S}$ be its positive-time reference test function. Then:
\begin{enumerate}[label=\textup{(\roman*)},leftmargin=2.25em]
\item
\[
\|E^{\mathrm{ren}}(\varphi_{*,S})\Omega_{\mathrm{loc}}\|^2=1;
\]
\item the centered Euclidean and Minkowski two-point functions of $E^{\mathrm{ren}}$ are therefore nonzero;
\item the spectral measure of the nonzero vector $E^{\mathrm{ren}}(\varphi_{*,S})\Omega_{\mathrm{loc}}$ under both $H_{\mathrm{loc}}$ and the Minkowski Hamiltonian $P^0$ is a nonzero finite positive measure supported in $[m_*,\infty)$.
\end{enumerate}
In particular, the reconstructed local gauge-invariant sharp-local theory is non-trivial in the load-bearing sense required later in the paper: it contains a local gauge-invariant field that is not a multiple of the identity and has nonzero vacuum-created spectral weight above the gap.
\end{theorem}

\proofsynopsis{The theorem packages the canonical scalar witness produced by the Lane~A normalization scheme with the OS/Wightman reconstruction and the vacuum-gap transfer. The core paper keeps the theorem-facing witness statement because it is part of the public endpoint claim; the detailed normalization and spectral-measure proof remain in Companion~III.}
\proofhome{Companion III, Section~3.}

\paragraph{Downstream Euclidean clustering.}
The dense-core Euclidean clustering statement in \cref{thm:main}(B) is a downstream corollary of the vacuum-gap theorem on the same sharp-local state: once Corollary~\ref{cor:continuum-vacuum-gap} is known, the standard OS semigroup estimate on a dense positive-time probe core gives exponential decay at rate $m_*$. Its full proof remains outside the numbered body packet because the gap itself has already been established in Section~\ref{sec:qe3-transport}.

\subsection{Faithful-Wilson universality and exact local-net endpoint}

\begin{theorem}[Faithful-Wilson universality hypotheses]\label{thm:faithful-wilson-hypotheses}
Let $G$ be compact connected simple and let $\rho_1,\rho_2$ be faithful finite-dimensional unitary Wilson representations. Form the finite continuum-normalized action family built from the two faithful Wilson plaquette actions and impose the same admissible gradient-flow renormalization prescription. Then the non-collapse, one-shot entrance, and ultraviolet-matching hypotheses required by the universality comparison theorems hold uniformly on this faithful-Wilson comparison family.

Equivalently, the finite faithful-Wilson comparison family meets the Route~1 universality hypotheses simultaneously for both faithful regularizations, each at its own entrance block scale, with common family-wise constants obtained by taking maxima over the two faithful choices.
\end{theorem}

\proofsynopsis{This theorem is the theorem-facing reduction of the finite faithful comparison problem: the auxiliary faithful Wilson actions form one common continuum-normalized family, and the packet-1 plus Route~1 ultraviolet/entrance constants close uniformly on that finite family. The core paper does not reproduce the comparison-family proof; it keeps only the hypothesis closure needed by the faithful-Wilson universality theorem.}
\proofhome{Companion III, Section~4.}

\begin{theorem}[Qualitative faithful-Wilson universality of the continuum local theory]\label{thm:faithful-wilson-universality}
Let $G$ be compact connected simple and let $\rho_1,\rho_2$ be faithful finite-dimensional unitary Wilson representations. Construct the corresponding continuum theories using the same admissible gradient-flow renormalization scheme and the same coherent normalization conditions. Then there exists a unique $*$-algebra isomorphism
\[
\Phi_{\rho_1\to\rho_2}:\Aloc(G,\rho_1)\xrightarrow{\ \cong\ }\Aloc(G,\rho_2)
\]
which, on every finite engineering cap, sends the Lane~A renormalized field attached to an exact-dimension quotient label $[P]$ and test function $f$ to the renormalized field with the same label $[P]$ and the same test function $f$. The induced map intertwines all finite-family Schwinger functions, the regionwise local nets, the OS vacuum reconstructions, and the reconstructed Wightman fields.

In particular, within the admissible gradient-flow renormalization class and after coherent normalization, the continuum local gauge-invariant sharp-local Yang--Mills theory and its vacuum gap depend only on $G$ and not on the faithful Wilson representation used to regularize the ultraviolet lattice theory.
\end{theorem}

\proofsynopsis{The comparison-family hypothesis closure of \cref{thm:faithful-wilson-hypotheses}, together with the Lane~A inductive-union construction and the field-correspondence theorem, identifies the two faithful ultraviolet regularizations at the continuum sharp-local level. The standard uniqueness clause of OS reconstruction then upgrades equality of Schwinger families to unitary equivalence of the reconstructed vacuum theories.}
\proofhome{Companion III, Section~4.}

\begin{corollary}[Quantitative ledger remains $\rho$-indexed]\label{cor:rho-indexed-ledger}
Fix a compact connected simple group $G$ and choose a faithful finite-dimensional unitary Wilson regularization $\rho$. Then \cref{thm:faithful-wilson-universality} identifies all faithful choices with the same continuum local theory $\Aloc(G)$, but the explicit Packet~1 / Route~1 certificate retained in this manuscript is still recorded through the $\rho$-indexed constants
\[
\beta_*(G,\rho),\qquad q^*(G,\rho),\qquad \mathrm{Cent}(G,\rho),\qquad m_{\mathrm{blk}}(G,\rho),\qquad m_*(G,\rho;u_{\mathrm{ren}}).
\]
The present manuscript does not claim that these explicit quantitative ledgers are uniform over all faithful $\rho$. Different faithful Wilson regularizations yield canonically isomorphic continuum local theories, but may certify different explicit lower bounds on their common vacuum gap.
\end{corollary}

\begin{proof}
The first statement is \cref{thm:faithful-wilson-universality}. The second follows because the explicit constants are still read from the chosen faithful ultraviolet regularization through the ultraviolet gate, one-shot entrance, and continuum transport theorems of Sections~\ref{sec:uv}, \ref{sec:route1}, and \ref{sec:qe3-transport}. No representation-uniform quantitative comparison theorem for those constants is proved in the present manuscript.
\end{proof}

\begin{theorem}[Exact local-net endpoint and exclusion of extended-support sector data]\label{thm:exact-local-net-endpoint}
Fix a compact connected simple gauge group $G$. Choose any faithful finite-dimensional unitary Wilson representation $\rho_{\mathrm{fid}}$ and let
\[
\Aloc(G,\rho_{\mathrm{fid}}),\qquad (\mathcal H_{\mathrm{loc}}(G,\rho_{\mathrm{fid}}),\Omega_{\mathrm{loc}}(G,\rho_{\mathrm{fid}}),U_{G,\rho_{\mathrm{fid}}})
\]
be the sharp-local Euclidean/Wightman theory produced from that faithful ultraviolet regularization. By \cref{thm:faithful-wilson-universality}, different faithful choices of $\rho_{\mathrm{fid}}$ yield canonically isomorphic theories, so we may write
\[
\Aloc(G),\qquad (\mathcal H_{\mathrm{loc}}(G),\Omega_{\mathrm{loc}}(G),U_G)
\]
for that canonical faithful-$\rho$-independent local theory. Then:
\begin{enumerate}[label=\textup{(\roman*)},leftmargin=2.25em]
\item by Corollary~\ref{cor:field-correspondence}, the local Wightman fields in this theory are precisely the renormalized gauge-invariant local differential polynomials in the curvature and its covariant derivatives represented by Lane~A;
\item if $\rho$ is any finite-dimensional unitary Wilson representation of $G$ and $G_{\mathrm{eff}}:=G/\ker(\rho)$, then the local observable net obtained from the ultraviolet pair $(G,\rho)$ is canonically identified with $\Aloc(G_{\mathrm{eff}})$; in particular, for faithful $\rho$ this is just the canonical theory $\Aloc(G)$;
\item line operators, surface defects, and other global-form-sensitive observables are not elements of the bounded-region local field algebra constructed here, because their localization regions are extended rather than bounded.
\end{enumerate}
Consequently the theorem-level object constructed in this manuscript is the canonical local gauge-invariant sharp-local net $\Aloc(G)$ together with its OS/Wightman vacuum representation, and no extended-support nonlocal sector data enter that theorem-level existence-and-gap claim.
\end{theorem}

\proofsynopsis{Item~(i) is the field-correspondence output of Corollary~\ref{cor:field-correspondence}. Item~(ii) combines faithful-Wilson universality with the effective-group descent lemmas of Section~\ref{sec:scope}. Item~(iii) is a localization statement: the bounded-region local field algebra constructed in the sharp-local theory does not contain observables whose support is extended. The core paper keeps precisely that theorem-level endpoint and omits later AQFT commentary.}
\proofhome{Companion III, Section~4.}

\begin{corollary}[Group-only restatement of the endpoint theorem]\label{cor:group-only-endpoint}
Theorem~\ref{thm:main} is a theorem about the local gauge-invariant Yang--Mills theory attached to each compact connected simple gauge group $G$. Any faithful finite-dimensional unitary Wilson representation may be used as an auxiliary ultraviolet regularization and, within the admissible gradient-flow renormalization class and after coherent normalization, yields the same continuum local theory $\Aloc(G)$ up to the canonical faithful-Wilson identification of Theorem~\ref{thm:faithful-wilson-universality}. The explicit Route~1 certificate constants retained in \cref{thm:main} remain indexed by the chosen faithful ultraviolet regularization, as recorded in Corollary~\ref{cor:rho-indexed-ledger}. A nonfaithful ultraviolet choice yields instead the local theory of the effective quotient $G_{\mathrm{eff}}:=G/\ker(\rho)$. Differences between global forms survive only in nonlocal sector data not asserted by \cref{thm:main}$;$ differences between faithful ultraviolet Wilson representations survive only through the present quantitative certificate, not through the continuum local theory itself.
\end{corollary}

\begin{proof}
This is immediate from \cref{thm:faithful-wilson-universality,thm:exact-local-net-endpoint} together with Corollary~\ref{cor:rho-indexed-ledger}.
\end{proof}

Together these results fix the theorem-level object proved in this paper: a local gauge-invariant sharp-local Yang--Mills theory whose vacuum Hamiltonian has a spectral gap and whose continuum local net is independent of the faithful ultraviolet Wilson regularization. What remains is only to assemble \cref{thm:main} from the packet stack already stated.

\paragraph{Endpoint claim hygiene.}
Within the admissible gradient-flow renormalization class and after coherent normalization, the endpoint packet proves exactly one theorem-level object: the canonical bounded-region local gauge-invariant sharp-local net $\Aloc(G)$ together with its OS/Wightman vacuum representation and vacuum-sector mass gap. The faithful-Wilson universality theorem is qualitative only; Corollary~\ref{cor:rho-indexed-ledger} records separately that the explicit Packet~1 / Route~1 quantitative ledger remains $\rho$-indexed. Extended-support line and surface sectors remain outside the theorem-level claim by Theorem~\ref{thm:exact-local-net-endpoint}.

